% This file provides examples of some useful macros for typesetting
% dissertations.  None of the macros defined here are necessary beyond
% for the template documentation, so feel free to change, remove, and add
% your own definitions.
%
% We recommend that you define macros to separate the semantics
% of the things you write from how they are presented.  For example,
% you'll see definitions below for a macro \file{}: by using
% \file{} consistently in the text, we can change how filenames
% are typeset simply by changing the definition of \file{} in
% this file.
%
%% The following is a directive for TeXShop to indicate the main file
%%!TEX root = main.tex

\newcommand{\NA}{\textsc{n/a}}  % for "not applicable"
\newcommand{\eg}{e.g.,\ }   % proper form of examples (\eg a, b, c)
\newcommand{\ie}{i.e.,\ }   % proper form for that is (\ie a, b, c)
\newcommand{\etal}{\emph{et al}}

% Some useful macros for typesetting terms.
\newcommand{\file}[1]{\texttt{#1}}
\newcommand{\class}[1]{\texttt{#1}}
\newcommand{\latexpackage}[1]{\href{http://www.ctan.org/macros/latex/contrib/#1}{\texttt{#1}}}
\newcommand{\latexmiscpackage}[1]{\href{http://www.ctan.org/macros/latex/contrib/misc/#1.sty}{\texttt{#1}}}
\newcommand{\env}[1]{\texttt{#1}}
\newcommand{\BibTeX}{Bib\TeX}

% Define a command \doi{} to typeset a digital object identifier (DOI).
% Note: if the following definition raises an error, then you likely
% have an ancient version of url.sty.  Either find a more recent version
% (3.1 or later works fine) and simply copy it into this directory,  or
% comment out the following two lines and uncomment the third.
\DeclareUrlCommand\DOI{}
\newcommand{\doi}[1]{\href{http://dx.doi.org/#1}{\DOI{doi:#1}}}
%\newcommand{\doi}[1]{\href{http://dx.doi.org/#1}{doi:#1}}

% Useful macro to reference an online document with a hyperlink
% as well with the URL explicitly listed in a footnote
% #1: the URL
% #2: the anchoring text
\newcommand{\webref}[2]{\href{#1}{#2}\footnote{\url{#1}}}

% epigraph is a nice environment for typesetting quotations
\makeatletter
\newenvironment{epigraph}{%
    \begin{flushright}
    \begin{minipage}{\columnwidth-0.75in}
    \begin{flushright}
    \@ifundefined{singlespacing}{}{\singlespacing}%
    }{
    \end{flushright}
    \end{minipage}
    \end{flushright}}
\makeatother

% \FIXME{} is a useful macro for noting things needing to be changed.
% The following definition will also output a warning to the console
\newcommand{\FIXME}[1]{\typeout{**FIXME** #1}\textbf{[FIXME: #1]}}

% END


\newcommand{\SimPEG}{\textsc{SimPEG}\xspace}
\newcommand{\simpegEM}{\textsc{simpegEM}\xspace}

\newcommand{\Mesh}{\texttt{Mesh}\xspace}
\newcommand{\Survey}{\texttt{Survey}\xspace}
\newcommand{\DCSurvey}{\texttt{DCSurvey}\xspace}
\newcommand{\Problem}{\texttt{Problem}\xspace}
\newcommand{\DCProblem}{\texttt{DCProblem}\xspace}
\newcommand{\EMProblem}{\texttt{EMProblem}\xspace}
\newcommand{\Regularization}{\texttt{Regularization}\xspace}
\newcommand{\DataMisfit}{\texttt{DataMisfit}\xspace}
\newcommand{\Optimization}{\texttt{Optimization}\xspace}
\newcommand{\InvProblem}{\texttt{InvProblem}\xspace}
\newcommand{\Inversion}{\texttt{Inversion}\xspace}

\newcommand{\Mapping}{\texttt{Mapping}\xspace}
\newcommand{\ExpMap}{\texttt{ExpMap}\xspace}
\newcommand{\PropMap}{\texttt{PropMap}\xspace}
\newcommand{\Sources}{\texttt{Sources}\xspace}
\newcommand{\Receivers}{\texttt{Receivers}\xspace}
\newcommand{\Fields}{\texttt{Fields}\xspace}

\newcommand{\TDEM}{\texttt{TDEM}\xspace}
\newcommand{\FDEM}{\texttt{FDEM}\xspace}
\newcommand{\NSEM}{\texttt{NSEM}\xspace}
\newcommand{\BaseEM}{\texttt{BaseEM}\xspace}

\newcommand{\Exp}{\texttt{Exp}\xspace}
\newcommand{\InjectActiveCells}{\texttt{InjectActiveCells}\xspace}

\newcommand{\dobs}{\mathbf{d}_\text{obs}}
\newcommand{\m}{\mathbf{m}}
\newcommand{\du}{\mathbf{u}}
\newcommand{\mref}{\mathbf{m}_\text{ref}}
\newcommand{\dpred}{\mathbf{d}_\text{pred}}
\newcommand{\meshI}{\emph{mesh}$_I$\xspace}
\newcommand{\meshF}{\emph{mesh}$_F$\xspace}
\newcommand{\Wd}{\mathbf{W}_\text{d}}
\newcommand{\Wm}{\mathbf{W}_\text{m}}
\newcommand{\sm}{\mathbf{s}_m}
\newcommand{\se}{\mathbf{s}_e}

\newcommand{\deriv}[2]{\frac{\partial #1}{\partial #2}}

\newcommand{\minimize}[1]{\mathop{\hbox{minimize}}_{#1}}

\renewcommand{\div}{\nabla\cdot\,}
% \newcommand{\grad}{\vec \nabla}
% \newcommand{\curl}{{\vec \nabla}\times\,}
\newcommand {\J}{{\vec J}}
\renewcommand{\H}{{\vec H}}
\newcommand {\E}{{\vec E}}
\renewcommand {\S}{{\vec{S}}}
\newcommand{\dcurl}{{\mathbf C}}
\newcommand{\dgrad}{{\mathbf G}}
\newcommand{\Acf}{{\mathbf A_c^f}}
\newcommand{\Ace}{{\mathbf A_c^e}}
\newcommand{\St}{{\mathbf \Sigma_\tau}}
\newcommand{\T}{{\mathbf T}}
\newcommand{\Tt}{{\mathbf T_\tau}}
\newcommand{\diag}[1]{\,{\sf diag}\left( #1 \right)}

\newcommand{\M}{{\mathbf M}}
\newcommand{\MfMui}{{\M^f_{\boldsymbol{\mu^{-1}}}}}
\newcommand{\MfRho}{{\M^f_{\boldsymbol{\rho}}}}
\newcommand{\MeSigma}{{\M^e_{\boldsymbol{\sigma}}}}
\newcommand{\MeMu}{{\M^e_{\boldsymbol{\mu}}}}
\newcommand{\mui}{\mu^{-1}}

\newcommand{\MeSigInf}{{\M^e_{\sigma_\infty}}}
\newcommand{\MeSigO}{{\M^e_{\sigma_0}}}
\newcommand{\Me}{{\M^e}}
\newcommand{\Mes}[1]{{\M^e_{#1}}}
\newcommand{\Mee}{{\M^e_e}}
\newcommand{\Mej}{{\M^e_j}}
\newcommand{\BigO}[1]{\mathcal{O}\bigl(#1\bigr)}
\newcommand{\bE}{\mathbf{E}}
\newcommand{\bH}{\mathbf{H}}
\newcommand{\B}{\vec{B}}
\newcommand{\D}{\vec{D}}
\renewcommand{\H}{\vec{H}}
\newcommand{\s}{\vec{s}}
\newcommand{\bfJ}{\bf{J}}
\newcommand{\vecm}{\vec m}
\renewcommand{\Re}{\mathsf{Re}}
\renewcommand{\Im}{\mathsf{Im}}
\renewcommand {\j}  { {\vec j} }
\newcommand {\h}  { {\vec h} }
\renewcommand {\b}  { {\vec b} }
\newcommand {\e}  { {\vec e} }
\newcommand {\db}  { {\mathbf {b}} }
\newcommand {\de}  { {\mathbf {e}} }
\renewcommand {\d}  { {\vec d} }
\renewcommand {\u}  { {\mathbf u} }
\newcommand{\I}{\vec{I}}

%================================Added by me
\newcolumntype{C}[1]{>{\centering\arraybackslash}p{#1}}
\newcolumntype{L}[1]{>{\raggedright\arraybackslash}p{#1}}
\newcolumntype{R}[1]{>{\raggedleft\arraybackslash}p{#1}}
\newcolumntype{J}[1]{>{\arraybackslash\justifying}p{#1}}
\newcolumntype{P}[1]{>{\centering\arraybackslash}p{#1}}

\newcommand{\hide}[1]{}
\newcommand{\replace}[2]{#2}
\newcommand{\replacewithnospace}[2]{\unskip#2}
\newenvironment{nitemize}{%
	\begin{itemize}[topsep=0pt,itemsep=2pt,parsep=0pt]%
	}{%
	\end{itemize}%
}

\newcommand*{\rom}[1]{\expandafter\romannumeral#1}
\newcommand{\comment}[1]{}
%%%%%%%%%%%%%%%%%%%%%%%%%%%%%%%%%%
\newcommand{\citet}[1]{
	% \citeauthor*{#1} \cite{#1}, \citeyear{#1} % IEEE Access
     \citeauthor{#1} \cite{#1}
}
\newcommand{\citex}[1]{
	% \citeauthor*{#1} \cite{#1}, \citeyear{#1} % IEEE Access
     \citeauthor{#1}, \citeyear{#1}, \cite{#1} %  Elsevier
}
%\newcommand{\GramaCheck}{}

\newcommand{\CheckRmv}[1]{#1}
% \newcommand{\figref}[1]{\hyperref[#1]{Figure~\ref{#1}}}%
% \newcommand{\tabref}[1]{\hyperref[#1]{Tabel~\ref{#1}}}%
\newcommand{\chapref}[1]{\hyperref[#1]{Chapter~\ref{#1}}}
% \newcommand{\secref}[1]{\hyperref[#1]{Section~\ref{#1}}}
% \renewcommand{\eqref}[1]{\hyperref[#1]{Equation~(\ref{#1})}}

\newcommand{\old}[1]{{\textcolor{purple}{old #1}}}
\newcommand{\red}[1]{{\textcolor{red}{#1}}}
% \newcommand{\cmm}[1]{{\textcolor{blue}{#1}}}
\newcommand{\gsh}[1]{{\textcolor{black}{#1}}} %
\newcommand{\kai}[1]{{\textcolor{awesome}{#1}}}
\newcommand{\ymh}[1]{{\textcolor{cyan}{#1}}}
\newcommand{\zxy}[1]{{\textcolor{green}{#1}}}
\newcommand{\myPara}[1]{\subsubsection{#1}}
\newcommand{\ourM}{{Res2Net}}
\newcommand{\sArt}{{state-of-the-art~}}
\newcommand{\tabSpace}{\vspace{6pt}}

\def\vs{\emph{vs.}}
\def\ie{\emph{i.e.,~}}
\def\eg{\emph{e.g.,~}}
\def\etc{\emph{etc}}
\def\etal{{\em et al.}}


\def\etal{{\em et al.}}
\newcommand{\tabFormat}{\centering \renewcommand{\arraystretch}{1.05}}
\usepackage[normalem]{ulem}

\DeclareGraphicsExtensions{.pdf,.jpg,.png}

\renewcommand{\tablename}{Table }

\newcommand{\figref}[1]{Figure.~\ref{#1}}
\newcommand{\tabref}[1]{Table.~\ref{#1}}
\newcommand{\equref}[1]{Equation.~(\ref{#1})}
\newcommand{\secref}[1]{Section.~\ref{#1}}
\newcommand{\reviewer}[1]{{\textcolor{black}{#1}}}
\newcommand{\mcols}[1]{\multicolumn{2}{c}{#1}}
\newcommand{\Cols}[1]{\multicolumn{2}{c|}{#1}}
\newcommand{\Rows}[1]{\multirow{2}{*}{#1}}
\newcommand{\colInf}{T50 & T100}
\newcommand{\ColInf}{\small{Flickr} & \small{Our} & \small{SHoG}}
\newcommand{\para}[1]{\noindent\textbf{#1.}}
\newcommand{\cmm}[1]{{\textcolor{blue}{#1}}}


\newcommand{\TSpace}{\vspace{-.05in}}
\newcommand{\BSpace}{\vspace{-.05in}}

\def\ie{\emph{i.e.~}}
\def\eg{\emph{e.g.~}}
\def\etal{{\em et al.~}}
\def\etc{{\em etc.~}}

\newcommand{\PutOverDes}{\small \put(13,0){Gau} \put(29,0){FT}\put(42,0){SEG}\put(56,0){\GRAB}\put(75,0){CB}\put(88,0){Ours}}

% ============================ By Saf1 ==========================================
\pdfminorversion=7 % To avoid the warning (PDF inclusion: found PDF version <1.7>, but at most version <1.5>)

%To create a hyperref for tables, figures, and equations
\newcommand{\reft}[1]{\hyperref[#1]{Table~\ref{#1}}}
\newcommand{\reff}[1]{\hyperref[#1]{Figure~\ref{#1}}}
\newcommand{\refe}[1]{\hyperref[#1]{Equation~(\ref{#1})}}
\newcommand{\refm}[1]{\hyperref[#1]{Matrix~(\ref{#1})}}
\newcommand{\refa}[1]{\hyperref[#1]{Algorithm~\ref{#1}}}
\newcommand{\refc}[1]{\hyperref[#1]{Chapter~\ref{#1}}}
\newcommand{\refs}[1]{\hyperref[#1]{Section~\ref{#1}}}
\newcommand{\refss}[1]{\hyperref[#1]{Subsection~\ref{#1}}}
\newcommand{\refap}[1]{\hyperref[#1]{Appendix~\ref{#1}}}
\newcommand{\refx}[2]{\hyperref[#1]{#1~\ref{#2}}} %for sections, pages..etc

\definecolor{richelectricblue}{rgb}{0.03, 0.57, 0.82}
\usepackage{hyperref}
\hypersetup{
    colorlinks=true,
    linkcolor=richelectricblue,
    urlcolor=richelectricblue,
    citecolor=richelectricblue,
    filecolor=richelectricblue
}